
\documentclass[11pt]{article}
\usepackage[cache=false]{minted}


\usepackage[letterpaper,margin=0.72in]{geometry}
\usepackage[T1]{fontenc}
\usepackage[utf8]{inputenc}
\usepackage{lmodern}
\usepackage{microtype}
\usepackage{xcolor}
\usepackage{hyperref}
\usepackage{bookmark}
\usepackage{fancyhdr}
\usepackage{titlesec}
\usepackage{enumitem}
\usepackage{longtable}
\usepackage{booktabs}
\usepackage{array}
\usepackage{calc}
\usepackage{amsmath,amssymb}
\usepackage{textcomp}
\usepackage{upquote}
% listings removed; minted is now the canonical backend
\usepackage[most]{tcolorbox}
\usepackage{etoolbox}

\definecolor{CABlue}{HTML}{1F4E79}
\definecolor{CALightBlue}{HTML}{EAF3F8}
\definecolor{CAGray}{HTML}{F7F7F7}
\definecolor{CADarkGray}{HTML}{444444}
\definecolor{CARule}{HTML}{B7C9D6}
\definecolor{CAAccent}{HTML}{355C7D}
\definecolor{CAGreen}{HTML}{2E6F40}
\definecolor{CAPurple}{HTML}{5B4B8A}

\hypersetup{
  colorlinks=true,
  linkcolor=CABlue,
  urlcolor=CABlue,
  citecolor=CABlue,
  pdftitle={Combinatorial Algorithms as C-like Pseudocode},
  pdfauthor={Jordan Suber / ChatGPT},
  pdfsubject={Algorithmic pseudocode handout based on Combinatorial Algorithms: For Computers and Calculators, Second Edition},
  pdfkeywords={combinatorial algorithms, Nijenhuis, Wilf, pseudocode, subsets, permutations, partitions, graphs, trees, network flow, backtracking}
}

\pagestyle{fancy}
\fancyhf{}
\lhead{\textit{Combinatorial Algorithms} Pseudocode}
\rhead{\thepage}
\renewcommand{\headrulewidth}{0.4pt}
\setlength{\headheight}{14pt}

\titleformat{\section}{\Large\bfseries\color{CABlue}}{}{0pt}{}
\titleformat{\subsection}{\large\bfseries\color{CABlue}}{}{0pt}{}
\titleformat{\subsubsection}{\normalsize\bfseries\color{CADarkGray}}{}{0pt}{}
\titlespacing*{\section}{0pt}{1.35em}{0.5em}
\titlespacing*{\subsection}{0pt}{0.95em}{0.35em}
\titlespacing*{\subsubsection}{0pt}{0.75em}{0.25em}

\setcounter{secnumdepth}{0}
\setcounter{tocdepth}{2}
\setlist[itemize]{leftmargin=1.4em,itemsep=0.2em,topsep=0.2em}
\setlist[enumerate]{leftmargin=1.6em,itemsep=0.25em,topsep=0.25em}
\setlength{\parindent}{0pt}
\setlength{\parskip}{0.52em}
\setlength{\LTpre}{0.35em}
\setlength{\LTpost}{0.35em}
\emergencystretch=4em
\sloppy


\newtcolorbox{scopebox}{
  enhanced,
  colback=CALightBlue,
  colframe=CABlue,
  boxrule=0.5pt,
  arc=2mm,
  left=1em,
  right=1em,
  top=0.7em,
  bottom=0.7em
}

\newtcolorbox{methodbox}{
  enhanced,
  colback=CAGray,
  colframe=CARule,
  boxrule=0.4pt,
  arc=2mm,
  left=0.9em,
  right=0.9em,
  top=0.6em,
  bottom=0.6em
}

\newtcolorbox{warningbox}{
  enhanced,
  colback=white,
  colframe=CAPurple,
  boxrule=0.45pt,
  arc=2mm,
  left=0.9em,
  right=0.9em,
  top=0.6em,
  bottom=0.6em
}

% Pandoc compatibility helpers
\providecommand{\tightlist}{\setlength{\itemsep}{0pt}\setlength{\parskip}{0pt}}
\providecommand{\passthrough}[1]{#1}
\AtBeginEnvironment{longtable}{\small}
\renewcommand{\arraystretch}{1.14}

\begin{document}
\hypersetup{pageanchor=false}
\pagenumbering{gobble}

\begin{titlepage}
  \centering
  \vspace*{1.7cm}
  {\Huge\bfseries\color{CABlue}Combinatorial Algorithms\par}
  \vspace{0.35cm}
  {\LARGE C-like Pseudocode Extraction\par}
  \vspace{0.25cm}
  {\large Based on Nijenhuis and Wilf, \textit{Combinatorial Algorithms: For Computers and Calculators}, Second Edition\par}
  \vspace{1.0cm}
  \begin{scopebox}
  This handout reorganizes the book's FORTRAN-oriented combinatorial routines into implementation-neutral C-like pseudocode. It is designed as a study, review, and implementation-planning artifact for generation, ranking, selection, graph, tree, polynomial, flow, and backtracking algorithms.
  \end{scopebox}
  \vspace{0.75cm}
  \begin{methodbox}
  \textbf{Scope:} systematic algorithm extraction from the uploaded scanned PDF, normalized into modern pseudocode.\\
  \textbf{Primary families:} subsets, $k$-subsets, compositions, permutations, integer partitions, set partitions, Young tableaux, heaps, cycles, spanning forests, chromatic polynomials, power series, network flow, permanents, posets, M\"obius functions, backtracking, labeled trees, random rooted trees, and minimum spanning trees.\\
  \textbf{Boundary:} the routines are rewritten as independent algorithm descriptions rather than transcribed FORTRAN source.
  \end{methodbox}
  \vfill
  {\large Prepared May 28, 2026\par}
\end{titlepage}

\hypersetup{pageanchor=true}
\pagenumbering{roman}
\tableofcontents
\clearpage
\pagenumbering{arabic}

\begin{warningbox}
\textbf{Implementation note.} The original book uses one-based FORTRAN arrays and compact stateful subroutines. This handout uses C-like control flow and mostly zero-based arrays. When implementing the algorithms, choose one indexing convention, add invariant checks, and validate counts against closed forms whenever possible.
\end{warningbox}

\section{Combinatorial Algorithms --- Systematic Algorithm Extraction}\label{combinatorial-algorithms-systematic-algorithm-extraction}

\textbf{Source basis.} This handout is based on the uploaded PDF of Albert Nijenhuis and Herbert S. Wilf, \emph{Combinatorial Algorithms: For Computers and Calculators}, Second Edition. The scanned source is organized around short FORTRAN subroutines, each chapter normally giving the mathematical basis, a formal algorithm, a flow chart, the FORTRAN routine, specifications, and sample output. This document rewrites the algorithmic content as independent C-like pseudocode. It is not a transcription of the book's FORTRAN source.

\textbf{Indexing convention.} The book uses FORTRAN-style 1-based arrays. The pseudocode below uses C-like 0-based arrays unless a ranking formula is clearer in 1-based notation. Intervals are half-open \texttt{[lo, hi)} unless explicitly marked inclusive.

\textbf{State convention for ``next'' routines.} The original routines commonly use a caller-maintained state flag meaning ``more to come.'' In this handout, each \texttt{next\_*} function returns \texttt{true} when it produced a valid next object and \texttt{false} when enumeration is exhausted. The first call is represented by an initialization function.

\textbf{Random convention.} Random generation routines below assume \texttt{uniform\_int(lo, hi)} returns an integer uniformly from the closed interval \texttt{[lo, hi]}, and \texttt{uniform01()} returns a real uniformly in \texttt{[0,1)}.

\begin{center}\rule{0.5\linewidth}{0.5pt}\end{center}

\subsection{\texorpdfstring{1. Subsets of an \texttt{n}-set}{1. Subsets of an n-set}}\label{subsets-of-an-n-set}

\subsubsection{1.1 Lexicographic next subset}\label{lexicographic-next-subset}

Represent a subset of \texttt{\{0,1,...,n-1\}} as a strictly increasing array \texttt{A[0..k-1]} of selected elements. Lexicographic order first changes the rightmost possible element, then resets the tail.

\begin{minted}{C}
bool next_subset_lex(int n, vector<int> *A) {
    int k = A->size();

    // Empty subset is followed by every one-element subset in this convention.
    if (k == 0) {
        A->push_back(0);
        return true;
    }

    // Try to extend the current subset by the smallest possible new element.
    if (A->back() < n - 1) {
        A->push_back(A->back() + 1);
        return true;
    }

    // Otherwise shorten until an element can be advanced.
    while (!A->empty() && A->back() == n - (int)A->size()) {
        A->pop_back();
    }
    if (A->empty()) return false;

    A->back() += 1;
    return true;
}
\end{minted}

\textbf{Output invariant.} \texttt{A} is always strictly increasing.\\
\textbf{Cost.} Amortized \texttt{O(1)} element changes; worst-case \texttt{O(n)} on rollover.

\subsubsection{1.2 Binary-counter subset enumeration}\label{binary-counter-subset-enumeration}

For compact implementation, use the subset mask as a binary counter.

\begin{minted}{C}
bool next_subset_mask(int n, uint64_t *mask) {
    if (*mask == ((1ULL << n) - 1ULL)) return false;
    *mask += 1;
    return true;
}
\end{minted}

\textbf{Use.} Exhaustive subset iteration when \texttt{n <= word\_bits}.\\
\textbf{Cost.} \texttt{O(1)} arithmetic per subset.

\subsubsection{1.3 Gray-code subset generation}\label{gray-code-subset-generation}

Successive subsets differ in exactly one element.

\begin{minted}{C}
uint64_t gray_subset(int rank) {
    return (uint64_t)rank ^ ((uint64_t)rank >> 1);
}

int changed_bit_between_gray(int rank) {
    // Change from gray(rank-1) to gray(rank) occurs at the lowest set bit of rank.
    return ctz(rank);
}

void enumerate_gray_subsets(int n, void visit(uint64_t mask)) {
    int total = 1 << n;
    for (int r = 0; r < total; r++) {
        visit(gray_subset(r));
    }
}
\end{minted}

\textbf{Use.} Iteration where incremental updates are cheaper than recomputation.

\begin{center}\rule{0.5\linewidth}{0.5pt}\end{center}

\subsection{\texorpdfstring{2. Random subset of an \texttt{n}-set}{2. Random subset of an n-set}}\label{random-subset-of-an-n-set}

\subsubsection{2.1 Independent Bernoulli subset}\label{independent-bernoulli-subset}

The unrestricted random subset is obtained by independently including each element with probability \texttt{1/2}.

\begin{minted}{C}
vector<int> random_subset(int n) {
    vector<int> A;
    for (int i = 0; i < n; i++) {
        if (uniform_int(0, 1) == 1) A.push_back(i);
    }
    return A;
}
\end{minted}

\textbf{Distribution.} Each of the \texttt{2\^n} subsets has probability \texttt{2\^\{-n\}}.

\subsubsection{\texorpdfstring{2.2 Random subset of prescribed probability \texttt{p}}{2.2 Random subset of prescribed probability p}}\label{random-subset-of-prescribed-probability-p}

\begin{minted}{C}
vector<int> random_subset_p(int n, double p) {
    vector<int> A;
    for (int i = 0; i < n; i++) {
        if (uniform01() < p) A.push_back(i);
    }
    return A;
}
\end{minted}

\begin{center}\rule{0.5\linewidth}{0.5pt}\end{center}

\subsection{\texorpdfstring{3. \texttt{k}-subsets of an \texttt{n}-set}{3. k-subsets of an n-set}}\label{k-subsets-of-an-n-set}

\subsubsection{\texorpdfstring{3.1 Lexicographic next \texttt{k}-subset}{3.1 Lexicographic next k-subset}}\label{lexicographic-next-k-subset}

\begin{minted}{C}
bool next_k_subset_lex(int n, int k, vector<int> *A) {
    // A must contain k increasing elements in [0,n).
    for (int i = k - 1; i >= 0; i--) {
        int max_at_i = n - k + i;
        if ((*A)[i] < max_at_i) {
            (*A)[i]++;
            for (int j = i + 1; j < k; j++) {
                (*A)[j] = (*A)[j - 1] + 1;
            }
            return true;
        }
    }
    return false;
}

vector<int> first_k_subset(int k) {
    vector<int> A(k);
    for (int i = 0; i < k; i++) A[i] = i;
    return A;
}
\end{minted}

\textbf{Cost.} Worst-case \texttt{O(k)}, amortized small.

\subsubsection{\texorpdfstring{3.2 Revolving-door next \texttt{k}-subset}{3.2 Revolving-door next k-subset}}\label{revolving-door-next-k-subset}

Revolving-door order changes exactly one element out and one element in. It is useful for incremental evaluation over combinations.

\begin{minted}{C}
struct RevolvingDoorState {
    int n, k;
    vector<int> A;      // current k-subset, sorted
    vector<int> dir;    // directions for mobile positions, values -1 or +1
    bool started;
};

void init_revolving(RevolvingDoorState *s, int n, int k) {
    s->n = n; s->k = k; s->started = false;
    s->A.resize(k);
    s->dir.assign(k, +1);
    for (int i = 0; i < k; i++) s->A[i] = i;
}

bool next_k_subset_revolving(RevolvingDoorState *s, int *out_removed, int *out_added) {
    if (!s->started) {
        s->started = true;
        *out_removed = -1; *out_added = -1;
        return true;
    }

    // Find a movable element whose one-step move preserves strict order and bounds.
    for (int i = s->k - 1; i >= 0; i--) {
        int candidate = s->A[i] + s->dir[i];
        bool within = (0 <= candidate && candidate < s->n);
        bool left_ok = (i == 0 || s->A[i - 1] < candidate);
        bool right_ok = (i == s->k - 1 || candidate < s->A[i + 1]);

        if (within && left_ok && right_ok) {
            *out_removed = s->A[i];
            s->A[i] = candidate;
            *out_added = candidate;
            return true;
        }
        s->dir[i] = -s->dir[i];
    }
    return false;
}
\end{minted}

\textbf{Note.} Exact revolving-door implementations differ in bookkeeping. The invariant is that successive combinations differ by a single exchange.

\begin{center}\rule{0.5\linewidth}{0.5pt}\end{center}

\subsection{\texorpdfstring{4. Random \texttt{k}-subset}{4. Random k-subset}}\label{random-k-subset}

\subsubsection{4.1 Sequential selection by remaining counts}\label{sequential-selection-by-remaining-counts}

At element \texttt{i}, choose it with probability \texttt{need / remaining}.

\begin{minted}{C}
vector<int> random_k_subset(int n, int k) {
    vector<int> A;
    int need = k;

    for (int i = 0; i < n; i++) {
        int remaining = n - i;
        if (need == 0) break;
        if (uniform_int(1, remaining) <= need) {
            A.push_back(i);
            need--;
        }
    }
    return A;
}
\end{minted}

\textbf{Distribution.} Uniform over all \texttt{C(n,k)} subsets.\\
\textbf{Cost.} \texttt{O(n)} time, \texttt{O(k)} output storage.

\subsubsection{4.2 Selection by combinatorial unranking}\label{selection-by-combinatorial-unranking}

\begin{minted}{C}
long long choose(int n, int k);  // assumed exact for relevant range

vector<int> unrank_k_subset(int n, int k, long long r) {
    vector<int> A;
    int x = 0;

    for (int pos = 0; pos < k; pos++) {
        for (int v = x; v < n; v++) {
            long long block = choose(n - v - 1, k - pos - 1);
            if (r < block) {
                A.push_back(v);
                x = v + 1;
                break;
            }
            r -= block;
        }
    }
    return A;
}

vector<int> random_k_subset_by_rank(int n, int k) {
    long long total = choose(n, k);
    return unrank_k_subset(n, k, uniform_int(0, total - 1));
}
\end{minted}

\begin{center}\rule{0.5\linewidth}{0.5pt}\end{center}

\subsection{\texorpdfstring{5. Compositions of \texttt{n} into \texttt{k} parts}{5. Compositions of n into k parts}}\label{compositions-of-n-into-k-parts}

A composition is an ordered vector of nonnegative or positive parts. The book's routines use ordered integer parts; both variants are useful.

\subsubsection{5.1 Next weak composition, lexicographic}\label{next-weak-composition-lexicographic}

Weak composition: \texttt{x[0] + ... + x[k-1] = n}, \texttt{x[i] >= 0}.

\begin{minted}{C}
vector<int> first_weak_composition(int n, int k) {
    vector<int> x(k, 0);
    x[0] = n;
    return x;
}

bool next_weak_composition(int n, int k, vector<int> *x) {
    // Move one unit from the leftmost possible nonlast part to the next part,
    // and collect all earlier mass back into part 0.
    for (int i = k - 2; i >= 0; i--) {
        if ((*x)[i] > 0) {
            (*x)[i]--;
            (*x)[i + 1]++;
            int sum_left = 0;
            for (int j = 0; j < i; j++) {
                sum_left += (*x)[j];
                (*x)[j] = 0;
            }
            (*x)[0] += sum_left;
            return true;
        }
    }
    return false;
}
\end{minted}

\subsubsection{5.2 Positive compositions via separators}\label{positive-compositions-via-separators}

Positive compositions of \texttt{n} into \texttt{k} positive parts correspond to choosing \texttt{k-1} dividers among \texttt{n-1} gaps.

\begin{minted}{C}
vector<int> composition_from_dividers(int n, vector<int> dividers) {
    // dividers are increasing positions in 1..n-1.
    vector<int> x;
    int prev = 0;
    for (int d : dividers) {
        x.push_back(d - prev);
        prev = d;
    }
    x.push_back(n - prev);
    return x;
}

bool next_positive_composition(int n, int k, vector<int> *x) {
    vector<int> d;
    int s = 0;
    for (int i = 0; i < k - 1; i++) {
        s += (*x)[i];
        d.push_back(s);
    }
    if (!next_k_subset_lex(n - 1, k - 1, &d)) return false;
    *x = composition_from_dividers(n, d);
    return true;
}
\end{minted}

\begin{center}\rule{0.5\linewidth}{0.5pt}\end{center}

\subsection{6. Random composition}\label{random-composition}

\subsubsection{6.1 Random positive composition}\label{random-positive-composition}

\begin{minted}{C}
vector<int> random_positive_composition(int n, int k) {
    vector<int> d = random_k_subset(n - 1, k - 1);
    for (int &v : d) v += 1;          // convert gap indices to divider positions
    sort(d.begin(), d.end());
    return composition_from_dividers(n, d);
}
\end{minted}

\subsubsection{6.2 Random weak composition}\label{random-weak-composition}

Weak compositions of \texttt{n} into \texttt{k} parts correspond to positive compositions of \texttt{n+k} into \texttt{k} parts after subtracting one from each part.

\begin{minted}{C}
vector<int> random_weak_composition(int n, int k) {
    vector<int> y = random_positive_composition(n + k, k);
    for (int &part : y) part -= 1;
    return y;
}
\end{minted}

\begin{center}\rule{0.5\linewidth}{0.5pt}\end{center}

\subsection{7. Permutations}\label{permutations}

\subsubsection{7.1 Lexicographic next permutation}\label{lexicographic-next-permutation}

\begin{minted}{C}
bool next_permutation_lex(vector<int> *p) {
    int n = p->size();
    int i = n - 2;
    while (i >= 0 && (*p)[i] >= (*p)[i + 1]) i--;
    if (i < 0) return false;

    int j = n - 1;
    while ((*p)[j] <= (*p)[i]) j--;

    swap((*p)[i], (*p)[j]);
    reverse(p->begin() + i + 1, p->end());
    return true;
}
\end{minted}

\subsubsection{7.2 Adjacent-transposition permutation generation}\label{adjacent-transposition-permutation-generation}

This Steinhaus-Johnson-Trotter style generator changes adjacent elements.

\begin{minted}{C}
bool next_permutation_adjacent(vector<int> *p, vector<int> *dir) {
    int n = p->size();
    int mobile = -1;
    int mobile_pos = -1;

    for (int i = 0; i < n; i++) {
        int j = i + (*dir)[(*p)[i]];
        if (0 <= j && j < n && (*p)[i] > (*p)[j]) {
            if ((*p)[i] > mobile) {
                mobile = (*p)[i];
                mobile_pos = i;
            }
        }
    }
    if (mobile_pos < 0) return false;

    int j = mobile_pos + (*dir)[mobile];
    swap((*p)[mobile_pos], (*p)[j]);

    for (int value = mobile + 1; value < n; value++) {
        (*dir)[value] = -(*dir)[value];
    }
    return true;
}
\end{minted}

\begin{center}\rule{0.5\linewidth}{0.5pt}\end{center}

\subsection{8. Random permutation}\label{random-permutation}

\subsubsection{8.1 Fisher-Yates shuffle}\label{fisher-yates-shuffle}

\begin{minted}{C}
vector<int> random_permutation(int n) {
    vector<int> p(n);
    for (int i = 0; i < n; i++) p[i] = i;

    for (int i = n - 1; i > 0; i--) {
        int j = uniform_int(0, i);
        swap(p[i], p[j]);
    }
    return p;
}
\end{minted}

\textbf{Distribution.} Uniform over all \passthrough{\mintinline"n!"} permutations.

\begin{center}\rule{0.5\linewidth}{0.5pt}\end{center}

\subsection{9. Integer partitions}\label{integer-partitions}

A partition of \texttt{n} is represented as a nonincreasing array of positive parts.

\subsubsection{9.1 Next integer partition}\label{next-integer-partition}

\begin{minted}{C}
vector<int> first_partition(int n) {
    return vector<int>{n};
}

bool next_partition(vector<int> *a) {
    // Generate partitions in reverse lexicographic refinement order.
    int k = a->size();
    int rem = 0;

    while (k > 0 && (*a)[k - 1] == 1) {
        rem += 1;
        k--;
    }
    if (k == 0) return false;

    (*a)[k - 1] -= 1;
    rem += 1;

    int x = (*a)[k - 1];
    a->resize(k);
    while (rem > x) {
        a->push_back(x);
        rem -= x;
    }
    if (rem > 0) a->push_back(rem);
    return true;
}
\end{minted}

\textbf{Invariant.} Parts remain positive and nonincreasing.

\subsubsection{9.2 Partition-count dynamic program}\label{partition-count-dynamic-program}

\begin{minted}{C}
vector<vector<long long>> partition_count_table(int n) {
    // P[s][m] = number of partitions of s with largest part at most m.
    vector<vector<long long>> P(n + 1, vector<long long>(n + 1, 0));
    for (int m = 0; m <= n; m++) P[0][m] = 1;

    for (int s = 1; s <= n; s++) {
        for (int m = 1; m <= n; m++) {
            P[s][m] = P[s][m - 1];
            if (s >= m) P[s][m] += P[s - m][m];
        }
    }
    return P;
}
\end{minted}

\begin{center}\rule{0.5\linewidth}{0.5pt}\end{center}

\subsection{10. Random integer partition}\label{random-integer-partition}

\subsubsection{10.1 Random partition by recursive unranking}\label{random-partition-by-recursive-unranking}

\begin{minted}{C}
vector<int> random_partition(int n) {
    vector<vector<long long>> P = partition_count_table(n);
    vector<int> parts;
    int remaining = n;
    int max_part = n;

    while (remaining > 0) {
        long long total = P[remaining][max_part];
        long long r = uniform_int(0, total - 1);

        int part = 1;
        for (int m = 1; m <= max_part; m++) {
            long long without_m_or_more = P[remaining][m - 1];
            long long with_part_m = (remaining >= m) ? P[remaining - m][m] : 0;
            if (r < without_m_or_more) continue;
            r -= without_m_or_more;
            if (r < with_part_m) { part = m; break; }
        }

        parts.push_back(part);
        remaining -= part;
        max_part = part;
    }
    sort(parts.rbegin(), parts.rend());
    return parts;
}
\end{minted}

\textbf{Distribution.} Uniform if the count table is exact and unranking intervals are handled without overflow.

\subsubsection{10.2 Next plane partition schema}\label{next-plane-partition-schema}

A plane partition can be represented as a matrix of nonnegative integers weakly decreasing along rows and columns.

\begin{minted}{C}
bool is_valid_plane_partition(Matrix<int> A) {
    for (int i = 0; i < A.rows; i++) {
        for (int j = 0; j < A.cols; j++) {
            if (j + 1 < A.cols && A[i][j] < A[i][j + 1]) return false;
            if (i + 1 < A.rows && A[i][j] < A[i + 1][j]) return false;
        }
    }
    return true;
}

bool next_plane_partition(Matrix<int> *A, int total_sum) {
    // Generic lexicographic backtracking successor.
    // Decrement the rightmost feasible cell and refill suffix greedily.
    vector<pair<int,int>> cells = row_major_cells(A->rows, A->cols);
    for (int t = cells.size() - 1; t >= 0; t--) {
        auto [i,j] = cells[t];
        if ((*A)[i][j] == 0) continue;
        (*A)[i][j]--;
        refill_remaining_cells_greedily(A, cells, t + 1, total_sum);
        if (sum(*A) == total_sum && is_valid_plane_partition(*A)) return true;
    }
    return false;
}
\end{minted}

\begin{center}\rule{0.5\linewidth}{0.5pt}\end{center}

\subsection{\texorpdfstring{11. Partitions of an \texttt{n}-set}{11. Partitions of an n-set}}\label{partitions-of-an-n-set}

Set partitions are represented by restricted-growth strings \texttt{a[0..n-1]}: \texttt{a[0]=0}, and \texttt{a[i] <= 1 + max(a[0..i-1])}.

\subsubsection{11.1 Next set partition}\label{next-set-partition}

\begin{minted}{C}
bool is_restricted_growth_prefix(vector<int> a, int end) {
    int max_seen = 0;
    if (a[0] != 0) return false;
    for (int i = 1; i < end; i++) {
        if (a[i] > max_seen + 1) return false;
        if (a[i] > max_seen) max_seen = a[i];
    }
    return true;
}

bool next_set_partition(vector<int> *a) {
    int n = a->size();
    vector<int> prefix_max(n, 0);
    for (int i = 1; i < n; i++) prefix_max[i] = max(prefix_max[i - 1], (*a)[i]);

    for (int i = n - 1; i >= 1; i--) {
        int max_allowed = prefix_max[i - 1] + 1;
        if ((*a)[i] < max_allowed) {
            (*a)[i]++;
            for (int j = i + 1; j < n; j++) (*a)[j] = 0;
            return true;
        }
    }
    return false;
}
\end{minted}

\textbf{Example.} \texttt{000}, \texttt{001}, \texttt{010}, \texttt{011}, \texttt{012} for \texttt{n=3}.

\subsubsection{11.2 Blocks from restricted-growth string}\label{blocks-from-restricted-growth-string}

\begin{minted}{C}
vector<vector<int>> blocks_from_rgs(vector<int> a) {
    int b = 1 + max_value(a);
    vector<vector<int>> blocks(b);
    for (int i = 0; i < a.size(); i++) blocks[a[i]].push_back(i);
    return blocks;
}
\end{minted}

\begin{center}\rule{0.5\linewidth}{0.5pt}\end{center}

\subsection{\texorpdfstring{12. Random partition of an \texttt{n}-set}{12. Random partition of an n-set}}\label{random-partition-of-an-n-set}

\subsubsection{12.1 Bell/Stirling count table}\label{bellstirling-count-table}

\begin{minted}{C}
vector<vector<long long>> stirling2_table(int n) {
    vector<vector<long long>> S(n + 1, vector<long long>(n + 1, 0));
    S[0][0] = 1;
    for (int i = 1; i <= n; i++) {
        for (int k = 1; k <= i; k++) {
            S[i][k] = k * S[i - 1][k] + S[i - 1][k - 1];
        }
    }
    return S;
}
\end{minted}

\subsubsection{12.2 Random set partition by number of blocks}\label{random-set-partition-by-number-of-blocks}

\begin{minted}{C}
vector<int> random_set_partition(int n) {
    auto S = stirling2_table(n);
    long long bell = 0;
    for (int k = 0; k <= n; k++) bell += S[n][k];

    long long r = uniform_int(0, bell - 1);
    int k = 0;
    while (r >= S[n][k]) { r -= S[n][k]; k++; }

    // Generate a random surjection onto k labeled blocks, then canonicalize labels.
    vector<int> a(n, -1);
    vector<int> block_size(k, 0);
    generate_random_surjection_rgs(n, k, &a, &block_size, S);
    canonicalize_rgs(&a);
    return a;
}
\end{minted}

\subsubsection{12.3 Recursive random restricted-growth string with fixed block count}\label{recursive-random-restricted-growth-string-with-fixed-block-count}

\begin{minted}{C}
void generate_random_surjection_rgs(int n, int k, vector<int> *a,
                                    vector<int> *block_size,
                                    vector<vector<long long>> S) {
    int used = 0;
    for (int i = 0; i < n; i++) {
        int remaining = n - i - 1;
        vector<pair<int,long long>> choices;

        for (int b = 0; b < used; b++) {
            long long count = count_completions(remaining, k, used, false, S);
            choices.push_back({b, count});
        }
        if (used < k) {
            long long count = count_completions(remaining, k, used + 1, true, S);
            choices.push_back({used, count});
        }

        int chosen = weighted_choice(choices);
        (*a)[i] = chosen;
        if (chosen == used) used++;
    }
}
\end{minted}

\begin{center}\rule{0.5\linewidth}{0.5pt}\end{center}

\subsection{13. General sequencing, ranking, and selection}\label{general-sequencing-ranking-and-selection}

The book abstracts many combinatorial families as rooted generation trees. Each object corresponds to a path; ranking and unranking count the number of leaves under skipped branches.

\subsubsection{13.1 Generic lexicographic unranking}\label{generic-lexicographic-unranking}

\begin{minted}{C}
typedef struct Node Node;

long long count_extensions(Node state);
vector<Move> legal_moves(Node state);
Node apply_move(Node state, Move m);
bool is_complete(Node state);
Object decode(Node state);

Object unrank_object(Node root, long long rank) {
    Node state = root;
    while (!is_complete(state)) {
        for (Move m : legal_moves(state)) {
            Node child = apply_move(state, m);
            long long block = count_extensions(child);
            if (rank < block) {
                state = child;
                break;
            }
            rank -= block;
        }
    }
    return decode(state);
}
\end{minted}

\subsubsection{13.2 Generic rank}\label{generic-rank}

\begin{minted}{C}
long long rank_object(Node root, Object x) {
    Node state = root;
    long long rank = 0;

    while (!is_complete(state)) {
        Move actual = move_used_by_object(state, x);
        for (Move m : legal_moves(state)) {
            if (m == actual) break;
            rank += count_extensions(apply_move(state, m));
        }
        state = apply_move(state, actual);
    }
    return rank;
}
\end{minted}

\subsubsection{13.3 Uniform random object by counts}\label{uniform-random-object-by-counts}

\begin{minted}{C}
Object random_object(Node root) {
    long long total = count_extensions(root);
    long long r = uniform_int(0, total - 1);
    return unrank_object(root, r);
}
\end{minted}

\begin{center}\rule{0.5\linewidth}{0.5pt}\end{center}

\subsection{14. Young tableaux}\label{young-tableaux}

A standard Young tableau of shape \texttt{lambda} fills cells with \texttt{1..n} increasing across rows and down columns.

\subsubsection{14.1 Hook-length count}\label{hook-length-count}

\begin{minted}{C}
long long hook_length_count(vector<int> lambda) {
    int n = sum(lambda);
    BigInt numerator = factorial(n);
    BigInt denom = 1;

    for (int r = 0; r < lambda.size(); r++) {
        for (int c = 0; c < lambda[r]; c++) {
            int right = lambda[r] - c - 1;
            int below = 0;
            for (int rr = r + 1; rr < lambda.size(); rr++)
                if (lambda[rr] > c) below++;
            denom *= (right + below + 1);
        }
    }
    return numerator / denom;
}
\end{minted}

\subsubsection{14.2 Recursive tableau generation}\label{recursive-tableau-generation}

\begin{minted}{C}
void generate_tableaux(vector<int> shape, Matrix<int> T, int value, void visit(Matrix<int>)) {
    if (value == 0) {
        visit(T);
        return;
    }

    for (Cell cell : removable_corners(shape)) {
        Matrix<int> U = T;
        U[cell.r][cell.c] = value;
        vector<int> next_shape = remove_cell(shape, cell);
        generate_tableaux(next_shape, U, value - 1, visit);
    }
}
\end{minted}

\subsubsection{14.3 Random tableau by hook walk / weighted corner removal}\label{random-tableau-by-hook-walk-weighted-corner-removal}

\begin{minted}{C}
Matrix<int> random_standard_tableau(vector<int> shape) {
    int n = sum(shape);
    Matrix<int> T = empty_tableau(shape);

    for (int value = n; value >= 1; value--) {
        vector<Cell> corners = removable_corners(shape);
        vector<pair<Cell,BigInt>> weights;
        for (Cell c : corners) {
            vector<int> reduced = remove_cell(shape, c);
            weights.push_back({c, hook_length_count(reduced)});
        }
        Cell chosen = weighted_choice(weights);
        T[chosen.r][chosen.c] = value;
        shape = remove_cell(shape, chosen);
    }
    return T;
}
\end{minted}

\begin{center}\rule{0.5\linewidth}{0.5pt}\end{center}

\section{Part 2 --- Combinatorial Structures}\label{part-2-combinatorial-structures}

\subsection{15. Sorting and heaps}\label{sorting-and-heaps}

\subsubsection{15.1 Heap construction}\label{heap-construction}

\begin{minted}{C}
void sift_down(vector<Key> *A, int start, int end) {
    int root = start;
    while (2 * root + 1 <= end) {
        int child = 2 * root + 1;
        int swap_idx = root;

        if ((*A)[swap_idx] < (*A)[child]) swap_idx = child;
        if (child + 1 <= end && (*A)[swap_idx] < (*A)[child + 1]) swap_idx = child + 1;
        if (swap_idx == root) return;

        swap((*A)[root], (*A)[swap_idx]);
        root = swap_idx;
    }
}

void heapify(vector<Key> *A) {
    int n = A->size();
    for (int start = (n - 2) / 2; start >= 0; start--) {
        sift_down(A, start, n - 1);
        if (start == 0) break;
    }
}
\end{minted}

\subsubsection{15.2 Heapsort}\label{heapsort}

\begin{minted}{C}
void heapsort(vector<Key> *A) {
    heapify(A);
    for (int end = A->size() - 1; end > 0; end--) {
        swap((*A)[0], (*A)[end]);
        sift_down(A, 0, end - 1);
    }
}
\end{minted}

\subsubsection{15.3 Extract heap maximum}\label{extract-heap-maximum}

\begin{minted}{C}
Key extract_heap_max(vector<Key> *H) {
    Key ans = (*H)[0];
    (*H)[0] = H->back();
    H->pop_back();
    if (!H->empty()) sift_down(H, 0, H->size() - 1);
    return ans;
}
\end{minted}

\begin{center}\rule{0.5\linewidth}{0.5pt}\end{center}

\subsection{16. Cycle structure of a permutation}\label{cycle-structure-of-a-permutation}

\subsubsection{16.1 Tag cycles}\label{tag-cycles}

\begin{minted}{C}
vector<int> cycle_tags(vector<int> p) {
    int n = p.size();
    vector<int> tag(n, -1);
    int c = 0;

    for (int i = 0; i < n; i++) {
        if (tag[i] != -1) continue;
        int v = i;
        while (tag[v] == -1) {
            tag[v] = c;
            v = p[v];
        }
        c++;
    }
    return tag;
}
\end{minted}

\subsubsection{16.2 Invert permutation}\label{invert-permutation}

\begin{minted}{C}
vector<int> invert_permutation(vector<int> p) {
    int n = p.size();
    vector<int> inv(n);
    for (int i = 0; i < n; i++) inv[p[i]] = i;
    return inv;
}
\end{minted}

\subsubsection{16.3 Cycle decomposition}\label{cycle-decomposition}

\begin{minted}{C}
vector<vector<int>> cycles_of_permutation(vector<int> p) {
    int n = p.size();
    vector<bool> seen(n, false);
    vector<vector<int>> cycles;

    for (int i = 0; i < n; i++) {
        if (seen[i]) continue;
        vector<int> cyc;
        int v = i;
        while (!seen[v]) {
            seen[v] = true;
            cyc.push_back(v);
            v = p[v];
        }
        cycles.push_back(cyc);
    }
    return cycles;
}
\end{minted}

\begin{center}\rule{0.5\linewidth}{0.5pt}\end{center}

\subsection{17. Renumbering rows and columns of an array}\label{renumbering-rows-and-columns-of-an-array}

The renumbering problem treats nonzero array positions as edges of a bipartite graph between row vertices and column vertices. Renumbering groups connected rows and columns.

\begin{minted}{C}
struct Renumbering {
    vector<int> row_component;
    vector<int> col_component;
    vector<int> row_order;
    vector<int> col_order;
};

Renumbering renumber_rows_cols(Matrix<int> A) {
    int R = A.rows, C = A.cols;
    vector<vector<int>> row_to_cols(R), col_to_rows(C);

    for (int i = 0; i < R; i++) {
        for (int j = 0; j < C; j++) {
            if (A[i][j] != 0) {
                row_to_cols[i].push_back(j);
                col_to_rows[j].push_back(i);
            }
        }
    }

    Renumbering out;
    out.row_component.assign(R, -1);
    out.col_component.assign(C, -1);
    int comp = 0;

    for (int r0 = 0; r0 < R; r0++) {
        if (out.row_component[r0] != -1) continue;
        queue<pair<bool,int>> q;
        q.push({true, r0});
        out.row_component[r0] = comp;

        while (!q.empty()) {
            auto [is_row, v] = q.front(); q.pop();
            if (is_row) {
                out.row_order.push_back(v);
                for (int c : row_to_cols[v]) if (out.col_component[c] == -1) {
                    out.col_component[c] = comp;
                    q.push({false, c});
                }
            } else {
                out.col_order.push_back(v);
                for (int r : col_to_rows[v]) if (out.row_component[r] == -1) {
                    out.row_component[r] = comp;
                    q.push({true, r});
                }
            }
        }
        comp++;
    }
    return out;
}
\end{minted}

\begin{center}\rule{0.5\linewidth}{0.5pt}\end{center}

\subsection{18. Spanning forest of a graph}\label{spanning-forest-of-a-graph}

\subsubsection{18.1 Depth-first spanning forest}\label{depth-first-spanning-forest}

\begin{minted}{C}
struct ForestResult {
    vector<int> parent;
    vector<int> component;
    vector<pair<int,int>> tree_edges;
};

void dfs_span(int u, int comp, Graph G, ForestResult *R) {
    R->component[u] = comp;
    for (int v : G.adj[u]) {
        if (R->component[v] == -1) {
            R->parent[v] = u;
            R->tree_edges.push_back({u, v});
            dfs_span(v, comp, G, R);
        }
    }
}

ForestResult spanning_forest(Graph G) {
    ForestResult R;
    int n = G.n;
    R.parent.assign(n, -1);
    R.component.assign(n, -1);

    int comp = 0;
    for (int u = 0; u < n; u++) {
        if (R.component[u] == -1) {
            dfs_span(u, comp, G, &R);
            comp++;
        }
    }
    return R;
}
\end{minted}

\subsubsection{18.2 Breadth-first spanning forest}\label{breadth-first-spanning-forest}

\begin{minted}{C}
ForestResult spanning_forest_bfs(Graph G) {
    ForestResult R;
    R.parent.assign(G.n, -1);
    R.component.assign(G.n, -1);
    int comp = 0;

    for (int s = 0; s < G.n; s++) {
        if (R.component[s] != -1) continue;
        queue<int> q;
        q.push(s);
        R.component[s] = comp;
        while (!q.empty()) {
            int u = q.front(); q.pop();
            for (int v : G.adj[u]) {
                if (R.component[v] == -1) {
                    R.component[v] = comp;
                    R.parent[v] = u;
                    R.tree_edges.push_back({u, v});
                    q.push(v);
                }
            }
        }
        comp++;
    }
    return R;
}
\end{minted}

\begin{center}\rule{0.5\linewidth}{0.5pt}\end{center}

\subsection{19. Newton forms of a polynomial}\label{newton-forms-of-a-polynomial}

\subsubsection{19.1 Evaluate polynomial in Newton form}\label{evaluate-polynomial-in-newton-form}

Newton form: \texttt{p(x)=c[0]+c[1](x-a[0])+...+c[n-1] prod\_\{j<n-1\}(x-a[j])}.

\begin{minted}{C}
double eval_newton(vector<double> c, vector<double> a, double x) {
    int n = c.size();
    double y = c[n - 1];
    for (int i = n - 2; i >= 0; i--) {
        y = y * (x - a[i]) + c[i];
    }
    return y;
}
\end{minted}

\subsubsection{19.2 Divided differences / Newton interpolation}\label{divided-differences-newton-interpolation}

\begin{minted}{C}
vector<double> newton_coefficients(vector<double> x, vector<double> y) {
    int n = x.size();
    vector<double> c = y;

    for (int j = 1; j < n; j++) {
        for (int i = n - 1; i >= j; i--) {
            c[i] = (c[i] - c[i - 1]) / (x[i] - x[i - j]);
        }
    }
    return c;
}
\end{minted}

\subsubsection{19.3 Convert Newton form to power basis}\label{convert-newton-form-to-power-basis}

\begin{minted}{C}
vector<double> newton_to_power(vector<double> c, vector<double> a) {
    int n = c.size();
    vector<double> p(1, c[n - 1]);

    for (int i = n - 2; i >= 0; i--) {
        vector<double> q(p.size() + 1, 0.0);
        for (int j = 0; j < p.size(); j++) {
            q[j]     += -a[i] * p[j];
            q[j + 1] +=  p[j];
        }
        q[0] += c[i];
        p = q;
    }
    return p;
}
\end{minted}

\subsubsection{19.4 Stirling transform}\label{stirling-transform}

\begin{minted}{C}
vector<long long> stirling_second_to_power(vector<long long> a) {
    // If f(x)=sum_k a[k] x^(underline k), convert to ordinary powers.
    int n = a.size();
    vector<vector<long long>> s(n, vector<long long>(n, 0));
    s[0][0] = 1;
    for (int i = 1; i < n; i++) {
        for (int k = 1; k <= i; k++) {
            s[i][k] = s[i - 1][k - 1] - (i - 1) * s[i - 1][k];
        }
    }

    vector<long long> p(n, 0);
    for (int k = 0; k < n; k++)
        for (int j = 0; j <= k; j++)
            p[j] += a[k] * s[k][j];
    return p;
}
\end{minted}

\begin{center}\rule{0.5\linewidth}{0.5pt}\end{center}

\subsection{20. Chromatic polynomial of a graph}\label{chromatic-polynomial-of-a-graph}

\subsubsection{20.1 Deletion-contraction recurrence}\label{deletion-contraction-recurrence}

\begin{minted}{C}
Polynomial chromatic_polynomial(Graph G) {
    if (G.has_loop()) return Polynomial::zero();
    if (G.edge_count() == 0) return Polynomial::x_power(G.n);

    Edge e = choose_edge(G);
    Graph G_delete = G;
    G_delete.remove_edge(e);

    Graph G_contract = contract_edge(G, e);

    return chromatic_polynomial(G_delete) - chromatic_polynomial(G_contract);
}
\end{minted}

\textbf{Identity.} \texttt{chi\_G(q) = chi\_\{G-e\}(q) - chi\_\{G/e\}(q)}.\\
\textbf{Use.} Exact chromatic polynomial for small and moderate graphs.

\subsubsection{20.2 Chromatic polynomial with memoization}\label{chromatic-polynomial-with-memoization}

\begin{minted}{C}
Polynomial chromatic_polynomial_memo(Graph G, Map<CanonicalGraph, Polynomial> *memo) {
    CanonicalGraph key = canonicalize(G);
    if (memo->contains(key)) return (*memo)[key];

    Polynomial ans;
    if (G.has_loop()) ans = Polynomial::zero();
    else if (G.edge_count() == 0) ans = Polynomial::x_power(G.n);
    else {
        Edge e = choose_edge(G);
        ans = chromatic_polynomial_memo(delete_edge(G, e), memo)
            - chromatic_polynomial_memo(contract_edge(G, e), memo);
    }
    (*memo)[key] = ans;
    return ans;
}
\end{minted}

\begin{center}\rule{0.5\linewidth}{0.5pt}\end{center}

\subsection{21. Power series operations}\label{power-series-operations}

Let series be arrays \texttt{a[0..n]} representing \texttt{sum a[i] x\^i}.

\subsubsection{21.1 Product of power series}\label{product-of-power-series}

\begin{minted}{C}
vector<double> series_mul(vector<double> a, vector<double> b, int n) {
    vector<double> c(n + 1, 0.0);
    for (int i = 0; i <= n; i++) {
        for (int j = 0; j <= n - i; j++) {
            c[i + j] += a[i] * b[j];
        }
    }
    return c;
}
\end{minted}

\subsubsection{21.2 Reciprocal of a power series}\label{reciprocal-of-a-power-series}

\begin{minted}{C}
vector<double> series_recip(vector<double> a, int n) {
    vector<double> b(n + 1, 0.0);
    b[0] = 1.0 / a[0];
    for (int k = 1; k <= n; k++) {
        double s = 0.0;
        for (int i = 1; i <= k; i++) s += a[i] * b[k - i];
        b[k] = -s / a[0];
    }
    return b;
}
\end{minted}

\subsubsection{21.3 Quotient of power series}\label{quotient-of-power-series}

\begin{minted}{C}
vector<double> series_div(vector<double> a, vector<double> b, int n) {
    return series_mul(a, series_recip(b, n), n);
}
\end{minted}

\subsubsection{21.4 Derivative and integral}\label{derivative-and-integral}

\begin{minted}{C}
vector<double> series_derivative(vector<double> a) {
    int n = a.size() - 1;
    vector<double> d(max(0, n), 0.0);
    for (int i = 1; i <= n; i++) d[i - 1] = i * a[i];
    return d;
}

vector<double> series_integral(vector<double> a, double constant) {
    int n = a.size();
    vector<double> I(n + 1, 0.0);
    I[0] = constant;
    for (int i = 0; i < n; i++) I[i + 1] = a[i] / (i + 1);
    return I;
}
\end{minted}

\subsubsection{21.5 Logarithm and exponential}\label{logarithm-and-exponential}

\begin{minted}{C}
vector<double> series_log(vector<double> a, int n) {
    // Requires a[0] = 1.
    vector<double> q = series_div(series_derivative(a), a, n - 1);
    return series_integral(q, 0.0);
}

vector<double> series_exp(vector<double> a, int n) {
    // Requires a[0] = 0. Solve b' = a' b and b[0]=1.
    vector<double> b(n + 1, 0.0);
    b[0] = 1.0;
    for (int k = 1; k <= n; k++) {
        double s = 0.0;
        for (int i = 1; i <= k; i++) s += i * a[i] * b[k - i];
        b[k] = s / k;
    }
    return b;
}
\end{minted}

\begin{center}\rule{0.5\linewidth}{0.5pt}\end{center}

\subsection{22. Network flow}\label{network-flow}

\subsubsection{22.1 Residual network initialization}\label{residual-network-initialization}

\begin{minted}{C}
struct FlowEdge { int to, rev; double cap; };
struct FlowGraph { vector<vector<FlowEdge>> adj; };

void add_flow_edge(FlowGraph *G, int u, int v, double cap) {
    FlowEdge a = {v, (int)G->adj[v].size(), cap};
    FlowEdge b = {u, (int)G->adj[u].size(), 0.0};
    G->adj[u].push_back(a);
    G->adj[v].push_back(b);
}
\end{minted}

\subsubsection{22.2 Augmenting-path max flow}\label{augmenting-path-max-flow}

\begin{minted}{C}
double max_flow(FlowGraph *G, int s, int t) {
    double flow = 0.0;
    int n = G->adj.size();

    while (true) {
        vector<int> parent_v(n, -1), parent_e(n, -1);
        queue<int> q;
        q.push(s);
        parent_v[s] = s;

        while (!q.empty() && parent_v[t] == -1) {
            int u = q.front(); q.pop();
            for (int ei = 0; ei < G->adj[u].size(); ei++) {
                FlowEdge &e = G->adj[u][ei];
                if (e.cap > 0 && parent_v[e.to] == -1) {
                    parent_v[e.to] = u;
                    parent_e[e.to] = ei;
                    q.push(e.to);
                    if (e.to == t) break;
                }
            }
        }
        if (parent_v[t] == -1) break;

        double aug = INF;
        for (int v = t; v != s; v = parent_v[v]) {
            FlowEdge &e = G->adj[parent_v[v]][parent_e[v]];
            aug = min(aug, e.cap);
        }

        for (int v = t; v != s; v = parent_v[v]) {
            FlowEdge &e = G->adj[parent_v[v]][parent_e[v]];
            e.cap -= aug;
            G->adj[v][e.rev].cap += aug;
        }
        flow += aug;
    }
    return flow;
}
\end{minted}

\subsubsection{22.3 Min-cut extraction after max flow}\label{min-cut-extraction-after-max-flow}

\begin{minted}{C}
vector<bool> reachable_in_residual(FlowGraph G, int s) {
    vector<bool> seen(G.adj.size(), false);
    queue<int> q;
    q.push(s);
    seen[s] = true;

    while (!q.empty()) {
        int u = q.front(); q.pop();
        for (FlowEdge e : G.adj[u]) {
            if (e.cap > 0 && !seen[e.to]) {
                seen[e.to] = true;
                q.push(e.to);
            }
        }
    }
    return seen;
}
\end{minted}

\begin{center}\rule{0.5\linewidth}{0.5pt}\end{center}

\subsection{23. Permanent of a matrix}\label{permanent-of-a-matrix}

\subsubsection{23.1 Ryser inclusion-exclusion formula}\label{ryser-inclusion-exclusion-formula}

For an \texttt{n x n} matrix \texttt{A}, \texttt{per(A)=(-1)\^n sum\_\{S subset [n]\} (-1)\^|S| prod\_i sum\_\{j in S\} A[i][j]}.

\begin{minted}{C}
double permanent_ryser(Matrix<double> A) {
    int n = A.rows;
    double total = 0.0;

    for (uint64_t mask = 1; mask < (1ULL << n); mask++) {
        int bits = popcount(mask);
        double prod = 1.0;
        for (int i = 0; i < n; i++) {
            double row_sum = 0.0;
            for (int j = 0; j < n; j++)
                if (mask & (1ULL << j)) row_sum += A[i][j];
            prod *= row_sum;
        }
        total += ((n - bits) % 2 == 0 ? +prod : -prod);
    }
    return total;
}
\end{minted}

\subsubsection{23.2 Gray-code Ryser update}\label{gray-code-ryser-update}

\begin{minted}{C}
double permanent_gray_ryser(Matrix<double> A) {
    int n = A.rows;
    vector<double> row_sum(n, 0.0);
    double total = 0.0;
    uint64_t prev = 0;

    for (uint64_t r = 1; r < (1ULL << n); r++) {
        uint64_t mask = gray_subset(r);
        uint64_t diff = mask ^ prev;
        int j = ctz(diff);
        int sign_col = (mask & diff) ? +1 : -1;

        for (int i = 0; i < n; i++) row_sum[i] += sign_col * A[i][j];

        double prod = 1.0;
        for (int i = 0; i < n; i++) prod *= row_sum[i];

        int bits = popcount(mask);
        total += ((n - bits) % 2 == 0 ? +prod : -prod);
        prev = mask;
    }
    return total;
}
\end{minted}

\textbf{Cost.} \texttt{O(n 2\^n)} arithmetic with \texttt{O(n)} memory.

\begin{center}\rule{0.5\linewidth}{0.5pt}\end{center}

\subsection{24. Inverting a triangular array}\label{inverting-a-triangular-array}

For a unit lower-triangular matrix \texttt{A}, find \texttt{B=A\^\{-1\}} by recurrence.

\begin{minted}{C}
Matrix<double> invert_lower_triangular(Matrix<double> A) {
    int n = A.rows;
    Matrix<double> B(n, n, 0.0);

    for (int i = 0; i < n; i++) {
        B[i][i] = 1.0 / A[i][i];
        for (int j = i - 1; j >= 0; j--) {
            double s = 0.0;
            for (int k = j + 1; k <= i; k++) {
                s += A[i][k] * B[k][j];
            }
            B[i][j] = -s / A[i][i];
        }
    }
    return B;
}
\end{minted}

\textbf{Use.} Incidence algebra inversion, binomial/Stirling transforms, and triangular recurrence inversion.

\begin{center}\rule{0.5\linewidth}{0.5pt}\end{center}

\subsection{25. Triangular numbering in partially ordered sets}\label{triangular-numbering-in-partially-ordered-sets}

A finite poset can be topologically numbered so that \texttt{x < y} implies \texttt{number[x] < number[y]}. Such numbering supports triangular storage of incidence matrices.

\subsubsection{25.1 Topological numbering}\label{topological-numbering}

\begin{minted}{C}
vector<int> triangular_numbering(Poset P) {
    int n = P.n;
    vector<int> indeg(n, 0);
    for (int u = 0; u < n; u++)
        for (int v : P.covers[u]) indeg[v]++;

    priority_queue<int, vector<int>, greater<int>> q;
    for (int u = 0; u < n; u++) if (indeg[u] == 0) q.push(u);

    vector<int> number(n, -1);
    int next = 0;
    while (!q.empty()) {
        int u = q.top(); q.pop();
        number[u] = next++;
        for (int v : P.covers[u]) {
            if (--indeg[v] == 0) q.push(v);
        }
    }
    if (next != n) error("relation is not acyclic");
    return number;
}
\end{minted}

\subsubsection{25.2 Transitive closure from triangular order}\label{transitive-closure-from-triangular-order}

\begin{minted}{C}
Matrix<bool> transitive_closure_poset(Poset P, vector<int> order) {
    int n = P.n;
    Matrix<bool> reach(n, n, false);
    for (int u = 0; u < n; u++) reach[u][u] = true;
    for (int u = 0; u < n; u++)
        for (int v : P.covers[u]) reach[u][v] = true;

    for (int kk = 0; kk < n; kk++) {
        int k = order[kk];
        for (int i = 0; i < n; i++) if (reach[i][k])
            for (int j = 0; j < n; j++)
                reach[i][j] = reach[i][j] || reach[k][j];
    }
    return reach;
}
\end{minted}

\begin{center}\rule{0.5\linewidth}{0.5pt}\end{center}

\subsection{26. Mobius function of a poset}\label{muxf6bius-function-of-a-poset}

For a finite poset, the Mobius function satisfies \texttt{mu(x,x)=1} and \texttt{sum\_\{x<=z<=y\} mu(x,z)=0} for \texttt{x<y}.

\begin{minted}{C}
Matrix<long long> mobius_function(Matrix<bool> leq, vector<int> topo) {
    int n = leq.rows;
    Matrix<long long> mu(n, n, 0);

    for (int x : topo) {
        mu[x][x] = 1;
        for (int y : topo) {
            if (x == y || !leq[x][y]) continue;
            long long s = 0;
            for (int z : topo) {
                if (z == y) continue;
                if (leq[x][z] && leq[z][y]) s += mu[x][z];
            }
            mu[x][y] = -s;
        }
    }
    return mu;
}
\end{minted}

\textbf{Use.} Inclusion-exclusion on posets and inversion in incidence algebras.

\begin{center}\rule{0.5\linewidth}{0.5pt}\end{center}

\subsection{27. Backtracking method}\label{backtracking-method}

\subsubsection{27.1 Generic backtracking kernel}\label{generic-backtracking-kernel}

\begin{minted}{C}
void backtrack(State s, void visit(Solution)) {
    if (is_solution(s)) {
        visit(extract_solution(s));
        return;
    }

    vector<Candidate> C = create_candidates(s);
    for (Candidate c : C) {
        if (!acceptable(s, c)) continue;
        apply(&s, c);
        backtrack(s, visit);
        undo(&s, c);
    }
}
\end{minted}

\subsubsection{\texorpdfstring{27.2 \texttt{n}-queens}{27.2 n-queens}}\label{n-queens}

\begin{minted}{C}
void queens(int row, int n, vector<int> *col_at_row,
            vector<bool> *used_col, vector<bool> *diag1, vector<bool> *diag2,
            void visit(vector<int>)) {
    if (row == n) { visit(*col_at_row); return; }

    for (int c = 0; c < n; c++) {
        int d1 = row - c + n - 1;
        int d2 = row + c;
        if ((*used_col)[c] || (*diag1)[d1] || (*diag2)[d2]) continue;

        (*col_at_row)[row] = c;
        (*used_col)[c] = (*diag1)[d1] = (*diag2)[d2] = true;
        queens(row + 1, n, col_at_row, used_col, diag1, diag2, visit);
        (*used_col)[c] = (*diag1)[d1] = (*diag2)[d2] = false;
    }
}
\end{minted}

\subsubsection{27.3 Hamiltonian circuits}\label{hamiltonian-circuits}

\begin{minted}{C}
void hamiltonian_cycles(Graph G, vector<int> *path, vector<bool> *used, void visit(vector<int>)) {
    int n = G.n;
    if (path->size() == n) {
        if (G.has_edge(path->back(), (*path)[0])) visit(*path);
        return;
    }

    int u = path->back();
    for (int v : G.adj[u]) {
        if ((*used)[v]) continue;
        (*used)[v] = true;
        path->push_back(v);
        hamiltonian_cycles(G, path, used, visit);
        path->pop_back();
        (*used)[v] = false;
    }
}
\end{minted}

\subsubsection{27.4 Euler circuits}\label{euler-circuits}

\begin{minted}{C}
vector<int> euler_circuit(Graph G, int start) {
    vector<int> circuit, stack;
    stack.push_back(start);

    while (!stack.empty()) {
        int u = stack.back();
        if (G.degree(u) > 0) {
            Edge e = G.remove_any_incident_edge(u);
            stack.push_back(e.other(u));
        } else {
            circuit.push_back(u);
            stack.pop_back();
        }
    }
    reverse(circuit.begin(), circuit.end());
    return circuit;
}
\end{minted}

\subsubsection{27.5 Spanning trees by backtracking over edges}\label{spanning-trees-by-backtracking-over-edges}

\begin{minted}{C}
void enumerate_spanning_trees(Graph G, int edge_index, DSU dsu,
                              vector<Edge> *chosen, void visit(vector<Edge>)) {
    if (chosen->size() == G.n - 1) {
        if (dsu.component_count() == 1) visit(*chosen);
        return;
    }
    if (edge_index == G.edges.size()) return;

    Edge e = G.edges[edge_index];

    if (dsu.find(e.u) != dsu.find(e.v)) {
        DSU d2 = dsu;
        d2.unite(e.u, e.v);
        chosen->push_back(e);
        enumerate_spanning_trees(G, edge_index + 1, d2, chosen, visit);
        chosen->pop_back();
    }

    enumerate_spanning_trees(G, edge_index + 1, dsu, chosen, visit);
}
\end{minted}

\begin{center}\rule{0.5\linewidth}{0.5pt}\end{center}

\subsection{28. Labeled trees}\label{labeled-trees}

\subsubsection{28.1 Prufer decode}\label{pruxfcfer-decode}

\begin{minted}{C}
vector<Edge> prufer_decode(vector<int> code) {
    int n = code.size() + 2;
    vector<int> degree(n, 1);
    for (int x : code) degree[x]++;

    min_heap<int> leaves;
    for (int i = 0; i < n; i++) if (degree[i] == 1) leaves.push(i);

    vector<Edge> edges;
    for (int x : code) {
        int leaf = leaves.pop_min();
        edges.push_back({leaf, x});
        if (--degree[x] == 1) leaves.push(x);
    }

    int a = leaves.pop_min();
    int b = leaves.pop_min();
    edges.push_back({a, b});
    return edges;
}
\end{minted}

\subsubsection{28.2 Prufer encode}\label{pruxfcfer-encode}

\begin{minted}{C}
vector<int> prufer_encode(Tree T) {
    int n = T.n;
    vector<int> degree = T.degrees();
    min_heap<int> leaves;
    for (int i = 0; i < n; i++) if (degree[i] == 1) leaves.push(i);

    vector<int> code;
    for (int step = 0; step < n - 2; step++) {
        int leaf = leaves.pop_min();
        int neighbor = T.only_live_neighbor(leaf, degree);
        code.push_back(neighbor);
        degree[leaf] = 0;
        if (--degree[neighbor] == 1) leaves.push(neighbor);
    }
    return code;
}
\end{minted}

\subsubsection{28.3 Enumerate labeled trees}\label{enumerate-labeled-trees}

\begin{minted}{C}
void enumerate_labeled_trees(int n, void visit(vector<Edge>)) {
    vector<int> code(n - 2, 0);
    while (true) {
        visit(prufer_decode(code));
        int i = n - 3;
        while (i >= 0 && code[i] == n - 1) i--;
        if (i < 0) break;
        code[i]++;
        for (int j = i + 1; j < n - 2; j++) code[j] = 0;
    }
}
\end{minted}

\begin{center}\rule{0.5\linewidth}{0.5pt}\end{center}

\subsection{29. Random unlabeled rooted trees}\label{random-unlabeled-rooted-trees}

Unlabeled rooted trees can be generated from integer partitions of subtree sizes. The exact implementation relies on counts by size and root degree.

\subsubsection{29.1 Count rooted trees by size}\label{count-rooted-trees-by-size}

\begin{minted}{C}
vector<BigInt> rooted_tree_counts(int N) {
    // Otter-style recurrence encoded as generating-function arithmetic.
    vector<BigInt> a(N + 1, 0);
    a[1] = 1;

    for (int n = 2; n <= N; n++) {
        BigInt s = 0;
        for (int k = 1; k <= n - 1; k++) {
            BigInt inner = 0;
            for (int d = 1; d <= k; d++) {
                if (k % d == 0) inner += d * a[d];
            }
            s += inner * a[n - k];
        }
        a[n] = s / (n - 1);
    }
    return a;
}
\end{minted}

\subsubsection{29.2 Random rooted tree by recursive partition of child subtrees}\label{random-rooted-tree-by-recursive-partition-of-child-subtrees}

\begin{minted}{C}
RootedTree random_unlabeled_rooted_tree(int n, CountOracle C) {
    if (n == 1) return single_node_tree();

    // Choose a multiset of child subtree sizes summing to n-1,
    // weighted by the number of unlabeled rooted trees of each size.
    vector<int> child_sizes = weighted_random_partition_of_root_children(n - 1, C);

    RootedTree root;
    for (int size : child_sizes) {
        root.children.push_back(random_unlabeled_rooted_tree(size, C));
    }
    sort_children_by_canonical_form(&root);
    return root;
}
\end{minted}

\textbf{Canonicalization.} Sorting child subtrees by canonical code avoids labeling artifacts.

\begin{center}\rule{0.5\linewidth}{0.5pt}\end{center}

\subsection{30. Tree of minimum length}\label{tree-of-minimum-length}

This is the minimum spanning tree problem for a weighted undirected graph.

\subsubsection{30.1 Kruskal minimum spanning tree}\label{kruskal-minimum-spanning-tree}

\begin{minted}{C}
vector<Edge> minimum_spanning_tree_kruskal(Graph G) {
    vector<Edge> edges = G.edges;
    sort(edges.begin(), edges.end(), by_weight_ascending);

    DSU dsu(G.n);
    vector<Edge> tree;

    for (Edge e : edges) {
        if (dsu.find(e.u) != dsu.find(e.v)) {
            dsu.unite(e.u, e.v);
            tree.push_back(e);
            if (tree.size() == G.n - 1) break;
        }
    }
    if (tree.size() != G.n - 1) error("graph is disconnected");
    return tree;
}
\end{minted}

\subsubsection{30.2 Prim minimum spanning tree}\label{prim-minimum-spanning-tree}

\begin{minted}{C}
vector<Edge> minimum_spanning_tree_prim(Graph G, int start) {
    vector<bool> in_tree(G.n, false);
    min_heap<pair<double, Edge>> pq;
    vector<Edge> tree;

    in_tree[start] = true;
    for (Edge e : G.incident(start)) pq.push({e.weight, e});

    while (!pq.empty() && tree.size() < G.n - 1) {
        auto [w, e] = pq.pop_min();
        int v = in_tree[e.u] ? e.v : e.u;
        if (in_tree[v]) continue;

        in_tree[v] = true;
        tree.push_back(e);
        for (Edge f : G.incident(v)) {
            int x = f.other(v);
            if (!in_tree[x]) pq.push({f.weight, f});
        }
    }
    if (tree.size() != G.n - 1) error("graph is disconnected");
    return tree;
}
\end{minted}

\begin{center}\rule{0.5\linewidth}{0.5pt}\end{center}

\section{Cross-Cutting Implementation Notes}\label{cross-cutting-implementation-notes}

\subsection{A. Overflow-safe counting}\label{a.-overflow-safe-counting}

Many ranking and random routines require exact combinatorial counts. Use arbitrary-precision integers or saturating arithmetic when counts are only used for comparisons.

\begin{minted}{C}
BigInt safe_choose(int n, int k) {
    if (k < 0 || k > n) return 0;
    k = min(k, n - k);
    BigInt ans = 1;
    for (int i = 1; i <= k; i++) {
        ans = ans * (n - k + i) / i;
    }
    return ans;
}
\end{minted}

\subsection{B. Canonical object output}\label{b.-canonical-object-output}

\begin{minted}{C}
void canonicalize_subset(vector<int> *A) { sort(A->begin(), A->end()); }

void canonicalize_rgs(vector<int> *a) {
    map<int,int> renumber;
    int next = 0;
    for (int &x : *a) {
        if (!renumber.contains(x)) renumber[x] = next++;
        x = renumber[x];
    }
}
\end{minted}

\subsection{C. Uniform weighted choice}\label{c.-uniform-weighted-choice}

\begin{minted}{C}
template<class T>
T weighted_choice(vector<pair<T, BigInt>> choices) {
    BigInt total = 0;
    for (auto [x,w] : choices) total += w;
    BigInt r = random_bigint(0, total - 1);
    for (auto [x,w] : choices) {
        if (r < w) return x;
        r -= w;
    }
    error("unreachable");
}
\end{minted}

\subsection{D. Exhaustive versus random algorithms}\label{d.-exhaustive-versus-random-algorithms}

The book repeatedly separates two categories:

\begin{enumerate}
\def\labelenumi{\arabic{enumi}.}
\tightlist
\item
  \textbf{Next-object generation.} Maintain compact state and return the next combinatorial object without recomputing from scratch.
\item
  \textbf{Random generation.} Use counting, recursive selection, or bijections so that every object in the target family is equally likely.
\end{enumerate}

That distinction is preserved throughout the pseudocode above.

\begin{center}\rule{0.5\linewidth}{0.5pt}\end{center}

\section{Algorithm Inventory by Chapter}\label{algorithm-inventory-by-chapter}

\begin{longtable}[]{@{}
  >{\raggedleft\arraybackslash}p{(\columnwidth - 4\tabcolsep) * \real{0.4000}}
  >{\raggedright\arraybackslash}p{(\columnwidth - 4\tabcolsep) * \real{0.3000}}
  >{\raggedright\arraybackslash}p{(\columnwidth - 4\tabcolsep) * \real{0.3000}}@{}}
\toprule\noalign{}
\begin{minipage}[b]{\linewidth}\raggedleft
Chapter
\end{minipage} & \begin{minipage}[b]{\linewidth}\raggedright
Book routine family
\end{minipage} & \begin{minipage}[b]{\linewidth}\raggedright
Pseudocode sections
\end{minipage} \\
\midrule\noalign{}
\endhead
\bottomrule\noalign{}
\endlastfoot
1 & \texttt{NEXSUB}, \texttt{LEXSUB} & Subset enumeration, binary and Gray-code variants \\
2 & \texttt{RANSUB} & Uniform random subset \\
3 & \texttt{NEXKSB}, \texttt{NXKSRD} & Lexicographic and revolving-door \texttt{k}-subsets \\
4 & \texttt{RANKSB} & Uniform random \texttt{k}-subset \\
5 & \texttt{NEXCOM} & Next composition \\
6 & \texttt{RANCOM} & Random composition \\
7 & \texttt{NEXPER} & Next permutation \\
8 & \texttt{RANPER} & Random permutation \\
9 & \texttt{NEXPAR} & Next integer partition \\
10 & \texttt{RANPAR}, \texttt{NEXT PLANE PARTITION} & Random integer partition; plane-partition schema \\
11 & \texttt{NEXEQU} & Next set partition / equivalence relation \\
12 & \texttt{RANEQU} & Random set partition \\
13 & \texttt{SELECT} & Generic sequencing, ranking, unranking, selection \\
14 & \texttt{NEXYTB}, \texttt{RANYTB} & Young tableaux generation and random selection \\
15 & \texttt{HPSORT}, \texttt{EXHEAP} & Heap operations and heapsort \\
16 & \texttt{CYCLES} & Cycle tags and permutation inversion \\
17 & \texttt{RENUMB} & Row/column renumbering as bipartite components \\
18 & \texttt{SPANFO} & Spanning forest and graph traversal \\
19 & \texttt{POLY} & Newton interpolation and polynomial transforms \\
20 & \texttt{CHROMP} & Chromatic polynomial \\
21 & \texttt{POWSER} & Formal power-series operations \\
22 & \texttt{NETFLO} & Network flow \\
23 & \texttt{PERMAN} & Permanent computation \\
24 & \texttt{INVERT} & Triangular-array inversion \\
25 & \texttt{TRIANG} & Triangular numbering for posets \\
26 & \texttt{MOBIUS} & Poset Mobius function \\
27 & \texttt{BACKTR} & Generic backtracking and examples \\
28 & \texttt{LBLTRE} & Labeled trees via Prufer codes \\
29 & \texttt{RANRUT} & Random unlabeled rooted trees \\
30 & \texttt{MINSPT} & Minimum spanning tree \\
\end{longtable}

\begin{center}\rule{0.5\linewidth}{0.5pt}\end{center}

\section{Practical Translation Checklist}\label{practical-translation-checklist}

\begin{itemize}
\tightlist
\item
  Replace all integer counts with \texttt{BigInt} for ranking, unranking, and random generation unless the input bounds are small.
\item
  Preserve one indexing convention throughout a real implementation; mixing 0-based pseudocode with 1-based formulas is a common source of off-by-one defects.
\item
  For random algorithms, test uniformity on small \texttt{n} by enumerating all outcomes and comparing empirical frequencies.
\item
  For next-object routines, test that enumeration count equals the appropriate closed form: \texttt{2\^n}, \texttt{C(n,k)}, \passthrough{\mintinline"n!"}, Bell number, partition number, or Cayley count \texttt{n\^\{n-2\}}.
\item
  For graph routines, normalize graph representation: directed versus undirected, loops allowed versus disallowed, and whether parallel edges are meaningful.
\item
  For recursive algorithms such as chromatic polynomial, permanent, random rooted trees, and backtracking, add memoization or pruning before using them on nontrivial instances.
\end{itemize}


\end{document}
